% ЗГЕНЕРОВАНО scripts/build_format_v2_chapter.py — не редагувати вручну.
% Джерело: books/ashtavakra-gita/poetychnyi/rozdil-3.md

\poemverse{3.1}{Пізнав ти Атман — незнищенний, єдиний. \\ То де тобі втіха з набутків, \\ атмазнавцю, мудрий?}
\poemverse{3.2}{Не знаєш Атмана — \\ і прикипаєш до речей у полі омани. \\ Не знаєш перламутру — \\ і жадібно хапаєш срібло, якого нема.}
\poemverse{3.3}{Пізнав: я — те, \\ у чому сяє весь цей світ, \\ як хвилі в морі. \\ Чого ж ти біжиш, наче вбогий?}
\poemverse{3.4}{Ти чув про Атман — \\ чисту свідомість, надзвичайно прекрасну. \\ Але хто надто привʼязаний до того, що під рукою, \\ той набирає бруду.}
\poemverse{3.5}{Мудрець знає Атман у всіх істотах \\ і всі істоти — в Атмані. \\ А своє досі тримається.}
\poemverse{3.6}{Став у вищій недвоїстості, \\ стоїш навіть у самій справі звільнення — \\ а бажання все ж веде тебе. \\ Безпорадний, мов у дитячій грі.}
\poemverse{3.7}{Зрозумів: знання було поганим другом. \\ І все ж, надзвичайно слабкий, \\ він прагнув би бажання \\ аж до кінця часу.}
\poemverse{3.8}{Хто відвернувся від цього світу й того, \\ хто розрізняє вічне й невічне, \\ хто хоче звільнення, — \\ той боїться саме звільнення.}
\poemverse{3.9}{Голублять його чи мучать — \\ він бачить тільки Атман. \\ Тому не радіє \\ і не гнівається.}
\poemverse{3.10}{Своє тіло, що рухається, \\ він бачить як тіло чуже. \\ То як його збентежить хвала \\ або ганьба?}
\poemverse{3.11}{Бачить: весь цей світ — тільки мая. \\ Охота до нього згасла. \\ То як злякається він, \\ коли смерть уже поруч, — \\ той, чий розум стійкий?}
\poemverse{3.12}{Розум без прагнень. \\ Великий дух — навіть у безнадії. \\ Вдоволений знанням Атмана. \\ З ким його й порівнювати?}
\poemverse{3.13}{Знає самою природою: \\ все, що видно, — аніщо. \\ Взяти це, відкинути те — \\ то що ж він бачить, \\ чий розум стійкий?}
\poemverse{3.14}{Каламуть усередині відкинуто. \\ Немає двоїстості, немає сподівань. \\ Все, що випадає само собою, \\ не стає ні болем, ні втіхою.}
